\documentclass[10pt]{article}
\usepackage[margin=20mm]{geometry}
\usepackage{amsmath,amssymb,fontspec,unicode-math,xurl,hyperref}
\IfFontExistsTF{DejaVu Serif}{\setmainfont{DejaVu Serif}}{\setmainfont{Latin Modern Roman}}
\setmathfont{latinmodern-math.otf}
\setlength{\parindent}{0pt}\setlength{\parskip}{6pt}\sloppy
\begin{document}
Abstract: The published LTR method, evaluation protocol, selected table arithmetic, and embedded limitations were inspected. No substantive invalid derivation or demonstrated core-result failure was established. The AGQA-Decomp compositional evaluation uses a supplied test graph, so it does not alone validate the complete generated-tree pipeline. Retrieval split and prediction-level reproducibility remain unresolved.\par
Author: AI-assisted technical review\par
Date:\par
Title: Technical Verification of Language-centric Tree Reasoning for Video Question-Answering\par
\# Paper identity and scope\par
This report concerns Zhaohe Liao et al., Divide and Conquer: Exploring Language-centric Tree Reasoning for Video Question-Answering, ICML 2025, PMLR 267, pp. 37511–37527, identifier I25-43530. The reviewed artifact is the staged published-proceedings PDF with embedded appendices. Page references use its printed pages 1–17, which coincide with the one-based PDF index. The source manifest labels the publication date 6 October 2025.\par
The review examined Sections 1–5, the method prompts in Sections 3.1–3.2, the evaluation protocol in Sections 4.1–4.3, Tables 1–5, the qualitative successes and failures, and Appendix A–C. Selected arithmetic and the interpretation of comparisons were checked. Models were not run and responses were not re-scored. No paper-specific implementation was present in the frozen sources; the acquisition record documents that availability limitation.\par
\# Checked reasoning\par
\#\# Operational method\par
The top-down stage retrieves question–subquestion examples from AGQA-Decomp and recursively asks an MLLM whether each subquestion is perceptual. The bottom-up stage answers leaves, combines their answers with the original video, and checks intermediate answers. These prompts define an inference procedure, not a formal soundness theorem. Because the video remains an input at internal nodes, a parent answer need not follow from textual leaf answers alone. This agrees with Section 3.2.2 (p. 5), Figure 3 and Section 5 (pp. 8–9), and the loss of accuracy when video is removed from bottom-up reasoning in Table 5 (p. 9).\par
\#\# Selected comparisons\par
The four matched backbone pairs in Table 1 give main-question accuracy gains of 3.8, 3.5, 3.5, and 3.7 percentage points and subquestion gains of 6.1, 6.0, 6.3, and 6.9 points. Independent decimal arithmetic reproduces these differences. Table 5 gives full-method gains of 2.2 and 2.3 points over the baseline on Causal-VidQA and NExT-QA. Removing answer verification lowers the full method by 0.5 and 0.7 points; removing video from bottom-up reasoning lowers it by 17.6 and 21.6 points. These checks support the direction of the displayed comparisons, without validating predictions or identifying a causal cognitive mechanism.\par
The qualitative failure case acknowledges that a missed book count propagates to an incorrect final answer. The paper is therefore not reasonably interpreted as claiming a correctness guarantee for the verification prompt. The sixteen-frame limitation for long videos is explicit in Appendix A.2 and C.1.\par
\# Scope clarification and unresolved evidence\par
\#\# Supplied graph in AGQA-Decomp\par
Section 4.3 (p. 6) expressly states that the DAG from the AGQA-Decomp test set is used for bottom-up reasoning when evaluating compositional consistency. Table 1 therefore measures reasoning with that supplied structure. It does not alone establish the complete generated-tree pipeline’s AGQA-Decomp accuracy. This is a disclosed protocol distinction, not demonstrated answer leakage. Other benchmarks and first-stage ablations supply separate evidence for the complete pipeline.\par
Section 3.1.1 (p. 4) describes embedding all non-leaf question–subquestion tuples in the retrieval database, but does not identify the retrieval split, exclusions, top-\(K\) value, or item-level retrieval records in the inspected materials. The supplied-graph protocol and retrieval database must be distinguished. The wording alone does not establish leakage. That question requires the exact retrieval inventory and evaluated question identifiers.\par
\#\# Interpretability and uncertainty\par
The tree exposes intermediate questions, answers, and explanations. That operational transparency is supported by the examples, but does not establish causal faithfulness of every explanation to the model’s internal computation. No quantitative faithfulness evaluation was available in this inspection.\par
Table captions distinguish copied results from reproduced backbone results. Matched pairs are the appropriate evidence for the framework’s incremental effect. Prediction records, judge configuration, and run-level uncertainty were unavailable. The compositional-score aggregation cannot be reconstructed from rounded entries alone: a harmonic mean of aggregate precision and recall need not equal an average of per-item harmonic means. The review does not label such an aggregation ambiguity a demonstrated arithmetic defect.\par
\# Assessment\par
No substantive invalid derivation or demonstrated core-result failure was established. The outcome is scope clarification, with No Issue Found in the declared inspected scope. The generated-tree AGQA-Decomp interpretation, retrieval split, exact compositional aggregation, and empirical reproducibility remain evidence-dependent. No correction is proposed. Later revisions and errata were not independently established in this inspection.\par
Issue | Assessment | Repair verification | Author adoption | Remaining work\par
Scope | No Issue Found | Not applicable | Unknown | Retrieval provenance and prediction-level reproducibility\par
Scientific adjudication is pending for every issue; human reviews remain empty. Author adoption has not been requested.\par
9 Z. Liao, J. Li, S. Sun, Q. Liu, F. Xiao, T. Li, Q. Zhang, G. Chen, L. Niu, C. Jiang, and L. Zhang, Divide and Conquer: Exploring Language-centric Tree Reasoning for Video Question-Answering, ICML, PMLR 267 (2025), 37511–37527. Published PDF including appendices. https://proceedings.mlr.press/v267/liao25d.html.\par
\# Evidence and execution\par
The parent conversation read the frozen task, checked actual source bytes against the manifest, and read the PDF with the installed `pdftotext -layout`. The direct PDF-reading connector failed; the source was unchanged. The original Table 1–2 page was rendered by the installed `pdftoppm` and visually inspected. An independently authored standard-library decimal-arithmetic script, `verification/0161/independent\_checks.py`, was executed. Its output is retained as `checks\_output.json`. This is finite arithmetic, not empirical reproduction.\par
Coverage and findings records distinguish inspected and unperformed work. Actual hashes and build state are retained in provenance. The parent AI performed this paper’s checks. No auxiliary scientific review or human review is asserted for this paper; adjudication remains pending and the human-review list is empty.\par
\end{document}
